\section{新朝不是一页插曲：从改名到道路失灵}
\label{sec:han-xin-depth}

王莽受禅时，长安的礼仪并不缺少文字。诏、册、符命、官号和印绶使改朝换代
看上去可以被排列成一套次序；真正难以排列的是这些文字离开都城之后会碰到的
土地关系、钱币习惯、地方官吏、边地道路和已经因饥荒而移动的人群。新朝之所以
值得单独阅读，不在于它是两汉之间的一段“失败”，而在于它把西汉长期积累的
问题同时拉到命令链上：名分能否代替信任，重新命名能否代替执行，官府账册能否
追上已经改变的家户与流动。

居摄时期，孺子婴仍保有刘氏的名义，王莽则控制解释名义的官署与人事。受禅仪式
将已经占住的位置说成天命的归属，不能证明地方或全部官僚出于自愿接受它。更
谨慎的说法是，王莽拥有足以把政治语言送出长安的中枢，却仍需依赖县廷、市场
和边郡把语言转成可执行的关系。土地和奴婢的改革、币制的更动、爵名官名的改造
都在这条链上接受检验；一项命名若要求所有人重新辨认交易、身份和文书，便会
使每一个中介同时面对服从、拖延、规避或借机牟利的选择。

\begin{event}{\eventref{WH-0009-XIN-USURPATION}{王莽受禅，改国号为新}}{九年}{新 · 长安}{}
受禅的队列、册命和国号改变，把宫廷已经发生的权力重组固定为公开仪式。它的
直接作用是令官员知道该用哪一种名义发令、领受印绶和安排祭祀；它无法替郡县
保证新的钱币会被接受，替边地保证旧有的盟友继续供给道路，或替正在流动的人群
提供土地。因而这不是“仪式无用”的证据，恰恰相反：仪式组织了谁可以用国家
的语言行动；只是国家语言还须在许多不由礼官控制的地点获得承载。

《汉书·王莽传》详细保存王莽的政治语言，却由东汉的复汉立场编成。新莽钱币、
官印和传世制度材料能显示某些名目确被制作和传达，也不能量出每一处市场的
接受程度。读者若将这两类材料放在一起，应当看见的是命令曾经抵达的痕迹和
它需要不断翻译的事实，而不是一幅“全国立即改革”或“全国从未执行”的图。
\end{event}

\subsection{钱币、田土与被改写的日常}

钱并不只是一种朝廷象征。商人要用它结算，家户要用它交换盐、布和粮，地方
官吏要用它折纳、赏赐或处罚。新朝多次改币，能够显示中枢试图重新规定价值和
流通；它也迫使持有旧钱、掌握铜料和负责市买的人反复判断何物可收、何物可拒。
后出的叙事常将市场紊乱写成改革本身的道德报应，不能从中计算全国物价或每个
家庭的损失。可把握的是，制度变动越需要基层不断辨认，越不能只靠都城的诏书
维持信用。

同样，土地与奴婢名目的调整触及的不是一张空白表格。土地已由家庭、豪右、佃作
和地方保护关系占有；人的身份又同债务、战乱、婚姻和劳役相连。禁止、限制或
重新分类的宣言，可能给某些人提供申诉的文字，也可能让官吏和强者取得新的
解释权。材料不允许我们把每项措施写成全国性胜败，更不允许把所有地方反应
合为一个“民怨”。它们要求我们看见执行不是政策的尾声：执行本身决定政策在
哪里被转译、在哪里被拖住、又在哪里生成新的冲突。

\subsection{荒年中的两种道路}

新朝后期的饥馑、洪水和兵乱，在不同区域并不具有同一成因和同一后果。关东的
流动人群、琅邪一带的赤眉、南方与西南的冲突、绿林与更始名号周围的武装，都
不能被一个“农民起义”的总称吞没。有人需要的是粮食和保护，有人掌握的是渡口、
乡里声望或一支队伍，有人则借旧汉、楚或新的名号把分散行动接成可谈判的力量。
朝廷若不能让粮、赈贷、任命和惩罚经由可信的地方中介到达，任何一个名号都会
变得比前一天更昂贵。

昆阳的战况、兵数和英雄言辞在《后汉书》与后来的叙事中极富戏剧性，不能拿来
重演每一处战场。它仍标明了中枢军政力量不能继续垄断局势的转折。更始政权进入
长安，也不是一座城被占领后全国便获得新主人：关中仓储、旧吏和居民要面对新的
征发与保护者，河北、陇右与巴蜀则继续拥有各自的军政道路。刘秀在河北取得的
优势，正来自他能把地方联盟、军队和一个重新可用的汉室名义逐段接合，而不是
因为“汉”字天然命令所有人服从。

\begin{sourcenote}[王莽改革为何不能只按成败读]
《汉书·王莽传》、帝纪和符命材料最清楚地保存中枢提出的名目；钱币、官印与
考古遗存证明部分制度形式进入物质世界；《后汉书》《后汉纪》则从复汉之后重排
崩解的年序。三者都不足以给出全国执行率、民意或财政收支。它们合在一起要求
读者区分：一项政策被宣布、被制作、被地方接收与改变地方关系，是四件不同的事。
\end{sourcenote}

\subsection{雒阳为什么能成为复接点}

刘秀在二十五年称帝，选择雒阳不是把旧都换成一座更合乎正统的城市。雒阳连着
关东与河北的道路，较容易接住支持者、粮运和能够迅速递送的文书；这也是它与
长安周边持续动荡相比的现实优势。赤眉进入长安、更始政权结束后，关中仍需经过
投降、安置和重新开通道路才可能回到新朝的联系中。陇右和巴蜀的收束更晚，说明
“复汉”并没有替每个地区省去自己重新选择官吏、军粮和保护者的过程。

新朝留下的不是一套与东汉无关的废墟。它使后来的雒阳朝廷明白，簿籍必须重新
取得人户，边地必须重新寻找中介，命令必须以不同速度经过不同道路。复汉初年的
度田和限制掠卖奴婢，正是这份教训在社会与行政层面的两种尝试：前者要让国家
重新看见土地和户口，后者要让战乱中被改写的人身关系能够提出申诉。两者都会
触动地方占有和吏员的判断，因此都不可能由一纸诏令完成。

\subsection{两座都城之间，谁还在办事}

长安在新朝崩解中并未立即失去一切行政痕迹，雒阳在复汉初年也并非从无到有地
创造秩序。两座城市之间的区别，首先在于哪一处能更稳定地接住军队、粮运、流民
和来自州郡的报告。更始政权进入长安时，需要利用城内的官署和仓储，也要面对
战争造成的空缺、物资紧张和不同军队的要求；赤眉进入关中时，同样不能只靠一
个名号解决供给。正史将政权更替写得紧凑，反而使读者更应留意：城邑被接收不
等于其周围的道路和家庭已经找到新的稳定关系。

刘秀把雒阳作为中枢，选择的是一套较能连接河北、河内、山东与南方交通的空间
条件。它不表示关中从此不重要，也不表示所有地方立刻向东看齐。隗嚣势力所在的
陇右、公孙述的巴蜀以及岭南的郡县，都拥有必须分别处理的山道、河流、地方精英
和驻军。复汉军队在不同区域获胜，只提供重新安排官吏和税粮的机会；若没有本地
愿意传送文书、供应军粮和承担司法的人，胜利便很难变成日常统治。

因此，新朝和复汉之间不宜用“混乱”与“统一”两个词截开。王莽后期的命令链
失灵，使人们在不同地方依靠不同保护者；雒阳朝廷的复接，则需要逐处承认这种
差异并寻找新的中介。建武度田的冲突、交趾的战事、河西以外重新试探的道路，
都是这项工作尚未完成的证据。东汉后来仍会重复面对同样的问题：一套制度能够
命名权力，却不能替代使命令获得粮食、信任与地方合作的关系。
